\section{Esercizio 1: Albero Binario, Generatori e Prototipi di Mappe}

\begin{definition}{Richieste Esercizio 1}
\textbf{1. Classe NodoBinario:} Rappresenta un nodo di un albero binario.
\begin{itemize}
    \item Costruttore: \texttt{val} (numero), \texttt{left} e \texttt{right} (opzionali, devono essere \texttt{NodoBinario} o \texttt{null}). Se \texttt{val} non è numero o i figli non sono validi, lanciare \texttt{TypeError}.
    \item Un nodo è \textbf{foglia} se ha entrambi i figli null.
\end{itemize}
\textbf{2. Generatore foglieConProfondita(radice):} 
\begin{itemize}
    \item Visita l'albero in \textbf{preordine}.
    \item Produce una sequenza di coppie \texttt{[profondita, valore]} per ogni nodo foglia (radice ha profondità 0).
\end{itemize}
\textbf{3. Estensione di Map (Prototipo):} 
\begin{itemize}
    \item Metodo \texttt{incrementa(key, amount)}: Se la chiave non c'è, la inserisce associata ad \texttt{amount}; se c'è, ne incrementa il valore. Ritorna la Map stessa.
\end{itemize}
\textbf{4. Metodo Statico in Map:}
\begin{itemize}
    \item \texttt{sommaFogliePerProfondita(radice)}: Usa il generatore per restituire una Map con chiave=profondità e valore=somma dei valori delle foglie a quella profondità.
\end{itemize}
\end{definition}


\subsection{Soluzione Codice}
\begin{jsbox}[title={Soluzione Esercizio 1}]
class NodoBinario {
  constructor(val, left = null, right = null) {
    if (typeof val !== "number") {
      throw new TypeError("val deve essere numero");
    }
    if (left !== null && !(left instanceof NodoBinario)) {
      throw new TypeError("left non valido");
    }
    if (right !== null && !(right instanceof NodoBinario)) {
      throw new TypeError("right non valido");
    }
    this.val = val;
    this.left = left;
    this.right = right;
  }

  // Un nodo è foglia se non ha figli
  get isLeaf() {
    return this.left === null && this.right === null;
  }
}

// Generatore per visitare le foglie in preordine
function* foglieConProfondita(nodo, profondita = 0) {
  if (nodo === null) return;
  
  // Se è foglia, produce la coppia [profondita, valore]
  if (nodo.isLeaf) {
    yield [profondita, nodo.val];
  } else {
    // Visita preordine: sinistra poi destra
    yield* foglieConProfondita(nodo.left, profondita + 1);
    yield* foglieConProfondita(nodo.right, profondita + 1);
  }
}

// Modifica del prototipo di Map
Map.prototype.incrementa = function(key, amount) {
  if (!this.has(key)) {
    this.set(key, amount);
  } else {
    this.set(key, this.get(key) + amount);
  }
  return this;
};

// Metodo statico aggiunto alla classe Map
Map.sommaFogliePerProfondita = function(radice) {
  let risultato = new Map();
  if (radice === null) return risultato;
  
  for (let [profondita, valore] of foglieConProfondita(radice)) {
    risultato.incrementa(profondita, valore);
  }
  return risultato;
};
\end{jsbox}

\section{Esercizio 2: Sistema di Ticketing, Ereditarietà ed Eccezioni}

\begin{definition}{Richieste Esercizio 2}
\textbf{1. Classe Ticket:}
\begin{itemize}
    \item Costruttore: \texttt{codice} (stringa non vuota), \texttt{priorita} (1-5 inclusi). Lancia \texttt{TypeError} se non validi.
    \item Proprietà private: \texttt{\#codice} (read-only), \texttt{\#chiuso} (read-only, inizialmente false), \texttt{\#log} (read-only array), \texttt{\#priorita} (read/write, lancia \texttt{PrioritaNonValida} se fuori 1-5).
    \item Metodi: \texttt{aggiungiNota(testo)} (aggiunge a log), \texttt{chiudi()} (imposta chiuso=true e logga "CHIUSURA").
    \item Se il ticket è chiuso, modifiche a priorità o log lanciano \texttt{ModificaTicketChiuso(codice)}.
\end{itemize}
\textbf{2. Classe TicketConSLA (extends Ticket):}
\begin{itemize}
    \item Aggiunge \texttt{\#tempoMassimo} (read-only) e \texttt{\#tempoTrascorso} (read/write, init 0).
    \item Se in scrittura \texttt{tempoTrascorso > tempoMassimo}, il ticket è \texttt{chiudi()} e lancia \texttt{SLASuperato(codice)}.
\end{itemize}
\textbf{3. Funzione incrementaTempo(tickets, delta):}
\begin{itemize}
    \item Prende array eterogeneo, incrementa il tempo solo alle istanze \texttt{TicketConSLA}.
    \item Non si interrompe all'errore, ma intercetta i \texttt{SLASuperato} e restituisce un array dei ticket che li hanno causati.
\end{itemize}
\end{definition}

\subsection{Soluzione Codice}
\begin{jsbox}[title={Soluzione Esercizio 2 - Classi ed Eccezioni}]
class PrioritaNonValida extends Error {}

class ModificaTicketChiuso extends Error {
  constructor(codice) {
    super("Ticket chiuso");
    this.codice = codice;
  }
}

class SLASuperato extends Error {
  constructor(codice) {
    super("SLA superato");
    this.codice = codice;
  }
}

class Ticket {
  #codice;
  #priorita;
  #chiuso;
  #log;

  constructor(codice, priorita) {
    if (typeof codice !== "string" || codice.length === 0) {
      throw new TypeError("codice non valido");
    }
    if (!Number.isInteger(priorita) || priorita < 1 || priorita > 5) {
      throw new TypeError("priorita non valida");
    }
    this.#codice = codice;
    this.#priorita = priorita;
    this.#chiuso = false;
    this.#log = [];
  }

  get codice() { return this.#codice; }
  get priorita() { return this.#priorita; }
  get chiuso() { return this.#chiuso; }
  // Restituisco una copia per proteggere l'array originale
  get log() { return [...this.#log]; }

  set priorita(p) {
    if (this.#chiuso) {
      throw new ModificaTicketChiuso(this.#codice);
    }
    if (!Number.isInteger(p) || p < 1 || p > 5) {
      throw new PrioritaNonValida();
    }
    this.#priorita = p;
    this.#log.push("Cambiata priorità a valore " + p);
  }

  aggiungiNota(testo) {
    if (this.#chiuso) {
      throw new ModificaTicketChiuso(this.#codice);
    }
    this.#log.push(testo);
  }

  chiudi() {
    if (!this.#chiuso) {
      this.#chiuso = true;
      this.#log.push("CHIUSURA");
    }
  }
}

class TicketConSLA extends Ticket {
  #tempoMassimo;
  #tempoTrascorso;

  constructor(codice, priorita, tempoMassimo) {
    super(codice, priorita);
    if (typeof tempoMassimo !== "number" || tempoMassimo <= 0) {
      throw new TypeError("tempoMassimo non valido");
    }
    this.#tempoMassimo = tempoMassimo;
    this.#tempoTrascorso = 0;
  }

  get tempoMassimo() { return this.#tempoMassimo; }
  get tempoTrascorso() { return this.#tempoTrascorso; }

  set tempoTrascorso(t) {
    if (typeof t !== "number" || t < 0) {
      throw new TypeError("tempoTrascorso non valido");
    }
    this.#tempoTrascorso = t;
    if (this.#tempoTrascorso > this.#tempoMassimo) {
      super.chiudi();
      throw new SLASuperato(this.codice);
    }
  }
}

function incrementaTempo(tickets, delta) {
  let risultati = [];
  if (!Array.isArray(tickets)) return risultati;
  
  for (let t of tickets) {
    if (t instanceof TicketConSLA) {
      try {
        t.tempoTrascorso = t.tempoTrascorso + delta;
      } catch (e) {
        if (e instanceof SLASuperato) {
          risultati.push(t);
        } else {
           throw e; // Rilancio se l'errore è inaspettato (non un SLA)
        }
      }
    }
  }
  return risultati;
}
\end{jsbox}

\section{Esercizio 3: Parcometro con Gestione dello Storico Globale}

\begin{definition}{Richieste Esercizio 3}
\textbf{1. Classe Movimento:}
Costruttore con: \texttt{tipo} ("R" o "S"), \texttt{euro} ($\ge 0$), \texttt{motivo} (stringa), \texttt{data} (istanza Date).
\\
\textbf{2. Classe Parcometro:}
\begin{itemize}
    \item Costruttore: \texttt{targa} (stringa non vuota), \texttt{credito} (numero $\ge 0$, default 0). Lancia TypeError se errati.
    \item Sola lettura per il credito disponibile.
    \item \texttt{ricarica(euro, motivo)} e \texttt{paga(euro, motivo)}: gestiscono il saldo e creano \texttt{Movimento}. Lanciano eccezioni sui tipi o \texttt{CreditoInsufficiente} se euro > credito (per \texttt{paga}).
    \item \texttt{storico(k=5)} restituisce copia sicura degli ultimi $k$ movimenti ordinati dal più recente.
    \item \textbf{Registro globale (\texttt{static \#registri = new Set()})}: Ogni operazione è una coppia \texttt{[p, m]} salvata nel set privato statico.
    \item \texttt{static registro()} restituisce un \texttt{new Set()} clonando il registro globale.
\end{itemize}
\end{definition}

\subsection{Soluzione Codice}
\begin{jsbox}[title={Soluzione Esercizio 3 - Metodi statici e Set}]
class CreditoInsufficiente extends Error {
  constructor(message = "Credito insufficiente") {
    super(message);
    this.name = "CreditoInsufficiente";
  }
}

class Movimento {
  constructor(t, e, m, d) {
    this.tipo = t;     // "R" ricarica oppure "S" spesa
    this.euro = e;     // e >= 0
    this.motivo = m;   // stringa
    this.data = d;     // Date
  }
}

class Parcometro {
  #targa;
  #movimenti = [];
  #creditoCorrente = 0;
  
  // Variabile statica privata per il registro globale (Set)
  static #registri = new Set();

  constructor(t, c = 0) {
    if (typeof t !== "string" || t === "") 
        throw new TypeError("Targa deve essere stringa non vuota!");
    if (typeof c !== "number" || c < 0) 
        throw new TypeError("Credito deve essere numero non negativo!");
        
    this.#targa = t;
    this.#creditoCorrente = c;
  }

  get credito() { return this.#creditoCorrente; }

  // Metodo statico per validazioni input
  static check(e, m) {
    if (typeof m !== "string") 
        throw new TypeError("Motivo deve essere stringa!");
    if (!Number.isFinite(e) || e <= 0) 
        throw new TypeError("Euro deve essere numero positivo!");
  }

  // Clona il set per non permettere alterazioni al registro interno
  static registro() {
    return new Set(this.#registri);
  }

  ricarica(euro, motivo) {
    Parcometro.check(euro, motivo);
    this.#creditoCorrente += euro;
    
    let mov = new Movimento("R", euro, motivo, new Date());
    this.#movimenti.push(mov);
    // Salva operazione nel registro globale statico
    Parcometro.#registri.add([this, mov]);
  }

  paga(euro, motivo) {
    Parcometro.check(euro, motivo);
    if (euro > this.#creditoCorrente) throw new CreditoInsufficiente();
    
    this.#creditoCorrente -= euro;
    
    let mov = new Movimento("S", euro, motivo, new Date());
    this.#movimenti.push(mov);
    Parcometro.#registri.add([this, mov]);
  }

  storico(k = 5) {
    // Restituisce array clonato: 
    // slice estrae gli ultimi k, map({...x}) clona gli oggetti Movimento
    // reverse() per averli dal più recente
    return [...this.#movimenti.slice(-k).map((x) => ({...x}))].reverse();
  }
}
\end{jsbox}
